\subsection{The ordinary Cohen--Macaulay comparison on odd dyadic nodes}

\begin{proposition}[Proved: the entire generic Swan kernel is invisible to the ordinary MCM comparison on odd nodes]
\label{C:odd-node-MCM-comparison-boundary}
Let $A$ be a characteristic-zero DVR with uniformizer $\pi$,
residue field $\kappa$ of characteristic two, and
$2=v\pi^e$, where $v\in A^\times$ and $e\geq1$.
Let $u\in A^\times$, $n\geq1$, and take either
\[
 R=\bigl(A[x,y]/(xy-u\pi^n)\bigr)_{(\pi,x,y)}
 \quad\text{or}\quad R=A[[x,y]]/(xy-u\pi^n).
\]
Put $B=R[C_2]$ and
\[
 K=\ker\bigl(R_R(C_2)\longrightarrow R_{R[1/2]}(C_2)\bigr),
 \qquad c:R_R(C_2)\longrightarrow G_0(B),             \tag{C.1.MB1}
\]
where $c$ is induced by inclusion into all finitely generated
$B$-modules.  Then:
\begin{enumerate}
\item The entire group
$\ker(G_0(B)\to G_0(B[1/\pi]))$ is cyclic of exponent
dividing $n$.  Tensoring with the sign representation acts
as the identity on this group.
\item The subgroup $c(K)$ is killed by both $2$ and $n$.
In particular, if $n$ is odd then $c(K)=0$.
\item If $\mathcal M$ is the exact category of finitely
generated $B$-modules that are maximal Cohen--Macaulay
over $R$, then $K_0(\mathcal M)\simeq G_0(B)$ under
inclusion.  Consequently the inclusion from $R$-free
$C_2$-lattices to $\mathcal M$ also kills $K$ when $n$
is odd.
\end{enumerate}
The latter inclusion need not be injective even for these
reduced normal domains: it is not injective on every
unramified node with odd $n\geq3$, by
Theorem~\ref{C:dyadic_normal_node_criterion}.
This is a boundary of the ordinary MCM comparison;
it does not assert that $K=0$ on any unresolved ramified node.
\end{proposition}
\begin{proof}
Replace $y$ by $z=y/u$, so that the equation is
$xz=\pi^n$.  The node domain and Cohen--Macaulay
arguments used in Theorem~\ref{C:dyadic_normal_node_criterion}
apply with $\pi$ as uniformizer.
In particular $R$ is a noetherian local domain of
dimension two, and $\pi,x,z$ are nonzero elements.
The ring $B$ is a commutative noetherian ring, finite
free of rank two over $R$.

\emph{Localization and the complete special-fiber calculation.}
Apply the $G_0$ localization sequence to the central
element $\pi$:
\[
 G_0(B/\pi B)\longrightarrow G_0(B)
       \longrightarrow G_0(B[1/\pi])\longrightarrow0. \tag{C.1.MB2}
\]
Precisely, the kernel category consists of finitely
generated $B$-modules killed by some power of $\pi$.
Each has a finite $\pi$-power filtration with factors
killed by $\pi$, and the quotient abelian category
is the category of finite $B[1/\pi]$-modules.
These are the noetherian and central-localization
hypotheses of the d\'evissage and localization statements
in \cite[II.6.3, II.6.3.3 and II.6.4.1]
{C-WeibelKbook2013}.

Write $C=R/\pi R$.  Thus
\[
 C=\bigl(\kappa[x,z]/(xz)\bigr)_{(x,z)}
 \quad\text{or}\quad C=\kappa[[x,z]]/(xz).
\]
If $t$ denotes the generator of $C_2$, then
\[
 B/\pi B=C[t]/((t-1)^2).
\]
The ideal generated by $t-1$ is nilpotent.
Its two-step filtration and nilpotent d\'evissage give
$G_0(B/\pi B)\simeq G_0(C)$, induced by regarding
a $C$-module as a representation with $t=1$
\cite[II.6.3.1]{C-WeibelKbook2013}.

Set
\[
 P=C/(x),\qquad Q=C/(z).
\]
Every finitely generated $C$-module has a finite prime
filtration.  The only prime ideals of this one-dimensional
local nodal ring are $(x)$, $(z)$ and $(x,z)$.
Consequently $G_0(C)$ is generated by $[P],[Q]$ and
the closed-point class $[\kappa]$.
There is no missing finite-length contribution:
the exact sequence
\[
 0\longrightarrow P\xrightarrow{z}P
       \longrightarrow\kappa\longrightarrow0
\]
gives $[\kappa]=0$.
Every finite-length module has a composition series
with factors $\kappa$, so every such class is zero.
The generic ranks on the two branches send $[P],[Q]$
to the two standard basis vectors of $\mathbb Z^2$;
hence $G_0(C)=\mathbb Z[P]\oplus\mathbb Z[Q]$.

We use the same letters for the images of these
classes in $G_0(B)$, always with trivial $C_2$-action.
Multiplication by $\pi$ on the trivial $B$-module
$R$ is injective.  Its cokernel is $C$, so $[C]=0$
in $G_0(B)$.
The exact sequence
\[
 0\longrightarrow Q\xrightarrow{\,\bar a\mapsto x\bar a\,}
 C\longrightarrow P\longrightarrow0
\]
therefore gives
\[
                         [P]+[Q]=0.                 \tag{C.1.MB3}
\]
Likewise multiplication by $x$ on the trivial
module $R$ is injective, and hence $[R/xR]=0$.
The ring $R/xR$ is the corresponding localization
or power-series ring over $A/(\pi^n)$ in the
variable $z$.
Its $\pi$-power filtration has exactly $n$ factors,
each isomorphic to $P$ as a $B$-module.
It follows that
\[
                              n[P]=0.                \tag{C.1.MB4}
\]
The surjectivity onto the kernel in (C.1.MB2),
together with (C.1.MB3)--(C.1.MB4), proves that
the entire generic $G_0$ kernel is cyclic with
exponent dividing $n$.  We do not assert that
its exponent is exactly $n$.

Tensoring with the free sign line is an exact
autoequivalence on finite $B$-modules.
On every module killed by $\pi$ it is literally
the identity action, because $-1=1$ in $C$.
The generation in (C.1.MB2) therefore proves
that sign twist acts trivially on this entire kernel.

\emph{The Swan augmentation relation.}
Let $s=[\mathrm{sgn}]$ in the Swan ring.
For any free rank-$d$ $C_2$-lattice $M$,
tensoring the regular augmentation sequence with
$M$ gives the admissible relation
\[
                         (1+s)[M]=d[R[C_2]].         \tag{C.1.MB5}
\]
The diagonal action on $R[C_2]\otimes_R M$ is
isomorphic to the induced regular action on the
underlying free $R$-module of rank $d$.
Every element $a\in K$ has underlying rank zero:
this follows already after passing to the fraction
field of the domain $R$.
By additivity, (C.1.MB5) gives $(1+s)a=0$.

Since $2=v\pi^e$, inverting two and inverting
$\pi$ give the same ring.
Thus $c(a)$ lies in the kernel from the first part.
There sign twist acts as the identity, so
\[
                         2c(a)=0,
                         \qquad nc(a)=0.             \tag{C.1.MB6}
\]
If $n$ is odd, Bezout's identity proves $c(a)=0$.
No assumption that the Swan kernel itself is
$2$-primary is used.

\emph{The MCM category really has this same group.}
The depth lemma over $R$ shows that $\mathcal M$
is closed under extensions and under kernels of
surjections between its objects.
Every finitely generated $B$-module $N$ has a
finite presentation by finite free $B$-modules.
In the resulting exact sequence
\[
 0\longrightarrow L\longrightarrow B^b
       \longrightarrow B^a\longrightarrow N\longrightarrow0,
\]
the second syzygy $L$ has $R$-depth at least two,
again by two applications of the depth lemma.
Here finite free $B$-modules have $R$-depth two
because $B$ is finite free over the Cohen--Macaulay
ring $R$.
Thus $L$ is in $\mathcal M$ and every finite
$B$-module has a finite $\mathcal M$-resolution.
The resolution theorem, with these closure and
finite-resolution hypotheses, gives
\[
                         K_0(\mathcal M)\simeq G_0(B)
                                                               \tag{C.1.MB7}
\]
under inclusion \cite[II.7.6]{C-WeibelKbook2013}.
The exact inclusion of $R$-free lattices into
$\mathcal M$ followed by (C.1.MB7) is exactly $c$.
This proves the third assertion.

\emph{A noninjective comparison on an explicit normal node.}
For example take
\[
 R_*=
 \bigl(\mathbb Z_{(2)}[x,y]/(xy+8)\bigr)_{(2,x,y)},
 \qquad
 T_* = \begin{pmatrix}3&x\\y&-3\end{pmatrix}.
\]
This is the explicit reduced normal domain at the
end of Theorem~\ref{C:dyadic_normal_node_criterion},
with $n=3$, $\pi=2$ and $u=-1$.
Its already proved mod-four primary parity detects
\[
              a_*=[T_*]-[\mathbf1]-[\mathrm{sgn}]\ne0.
\]
Indeed the specialization $x,y\mapsto2$ gives
the irreducible primary $z^2+z+1$ of length one,
whereas both constant sign lines have primary length
zero.  The explicit generic eigenbasis in that
theorem puts $a_*$ in $K$.
The present odd-node argument gives $c(a_*)=0$,
and therefore also zero in $K_0(\mathcal M)$.
The same earlier theorem supplies such a nonzero
generic class for every unramified odd $n\geq3$.
Thus the failure of injectivity occurs within
this same normal node family; it is not inferred
from a nonreduced example.
\end{proof}

\noindent\emph{Consequence for the remaining ramified case.}
On the exponent-three node the map from the generic
Swan kernel into ordinary $G_0$ or into ordinary
$K_0(\mathcal M)$ is zero for every ramification index.
Proving it injective on that kernel would therefore
already prove that the kernel vanishes.
A tame-quotient or matrix-factorization construction
that passes only through this ordinary MCM group
cannot supply a faithful obstruction without
additional exact-category information.
This statement leaves the ramified exponent-three
Swan comparison unresolved.
